% ============================================================================
%  SECTION 11 — Summary & references
% ============================================================================
\begin{frame}[t]{How the Four Theories Answer: The Takeaways}
  \small
  \begin{center}
  \renewcommand{\arraystretch}{1.5}
  \begin{tabular}{@{}p{0.24\textwidth}p{0.20\textwidth}p{0.48\textwidth}@{}}
    \toprule
    \textbf{Theory} & \textbf{Stability?} & \textbf{What it says} \\
    \midrule
    \textbf{\color{YaleBlue}Kinematic growth} & \emph{Imposed} &
      Either reaches the \emph{a-priori prescribed} homeostasis, or is unstable,
      nothing in between. Built in, not learned. \\
    \textbf{\color{YaleBlue}Constrained mixture (full)} & \emph{Predicted} &
      Adapts or dilates \textbf{depending on the insult}, the faithful
      reference. \\
    \textbf{\color{YaleBlue}Homogenized CMM} & \emph{Predicted} &
      A \textbf{similar time trajectory} to the full model, at a fraction of the
      cost. \\
    \textbf{\color{YaleBlue}Equilibrated CMM} & \emph{Instantly} &
      \textbf{Matches} the transient theories \emph{if} an equilibrium exists;
      \emph{no} root $\Rightarrow$ it does not match, unbounded growth. \\
    \bottomrule
  \end{tabular}
  \end{center}
\end{frame}

\begin{frame}[t]{Summary: What to Take Away}
  \textbf{The biology.} Living tissue continuously \textbf{turns over}; cells
  deposit new material \emph{pre-stretched} at $G^{\alpha}$, chasing a homeostatic
  mechanical set-point $\sigma_h$.
  \vfill
  \textbf{The four theories, one set-point, one question: does it adapt or run away?}
  \begin{itemize}\small\itemsep2pt
    \item \textbf{Kinematic growth}: prescribe $\bm{F}_g$; \emph{imposes} the
    set-point. Outcome is binary: reach the prescribed homeostasis, or go unstable.
    \item \textbf{Constrained mixture}: track every cohort; stability is
    \emph{predicted} and \textbf{depends on the insult} (mild $\to$ adapt, severe
    $\to$ runaway).
    \item \textbf{Homogenized}: one mean natural configuration per constituent;
    a \textbf{similar time trajectory} to the full model at $O(N)$ not $O(N^2)$.
    \item \textbf{Equilibrated}: solve for the end state directly; \textbf{matches}
    the transient theories \emph{iff} an equilibrium exists, and flags instability
    by finding no root.
  \end{itemize}
  \vfill
  \textbf{Where it goes.} These tools now drive \textbf{clinical prediction}
  (TEVGs), \textbf{systems-biology coupling}, and \textbf{fluid--solid--growth}.
\end{frame}

\begin{frame}[t,allowframebreaks]{Key references}
  \footnotesize
  \begin{itemize}\setlength{\itemsep}{0.35em}
    \item Rodriguez, Hoger, McCulloch. Stress-dependent finite growth in soft
    elastic tissues. \textit{J Biomech} (1994).
    \item Humphrey, Rajagopal. A constrained mixture model for growth and
    remodeling of soft tissues. \textit{Math Models Methods Appl Sci} (2002).
    \item Cyron, Wilson, Humphrey. Mechanobiological stability. \textit{J R Soc
    Interface} (2014).
    \item Cyron, Aydin, Humphrey. A homogenized constrained mixture (and
    mechanical analog) model. \textit{Biomech Model Mechanobiol} (2016).
    \item Latorre, Humphrey. A mechanobiologically equilibrated constrained
    mixture model. \textit{ZAMM} (2018).
    \item Latorre, Humphrey. Fast, rate-independent finite-element implementation
    of a 3D constrained mixture model. \textit{Comput Methods Appl Mech Eng} (2020).
    \item Humphrey. Constrained mixture models of soft tissue growth and
    remodeling. \textit{J Elasticity} (2021).
    \item Irons, Humphrey. Cell signaling model for arterial mechanobiology.
    \textit{PLoS Comput Biol} (2020).
    \item Irons, Latorre, Humphrey. From transcript to tissue: multiscale
    modeling. \textit{Ann Biomed Eng} (2021).
    \item Drews, Pepper, Best, \ldots, Humphrey, Breuer. Spontaneous reversal of
    stenosis in tissue-engineered vascular grafts. \textit{Sci Transl Med} (2020).
    \item Figueroa, Baek, Taylor, Humphrey. A computational framework for
    fluid--solid--growth modeling. \textit{Comput Methods Appl Mech Eng} (2009).
    \item Schwarz, Pfaller, \ldots, Marsden. A fluid--solid--growth solver for
    cardiovascular modeling. \textit{Comput Methods Appl Mech Eng} (2023).
    \item Pfaller, Latorre, \ldots, Marsden. FSGe: a fast, strongly coupled 3D
    fluid--solid--growth method. \textit{Comput Methods Appl Mech Eng} (2024).
    \item Humphrey, Eberth, Dye, Gleason. Fundamental role of axial stress in
    arterial mechanics. \textit{J Biomech} (2009).
    \item Ambrosi, Ben Amar, Cyron, \ldots, Humphrey, Kuhl. Growth and remodelling
    of living tissues. \textit{J R Soc Interface} (2019).
  \end{itemize}
  \framebreak
  \small
  \textbf{Code}
  \begin{itemize}
    \item Seminar code (this deck's figures): \texttt{github.com/YaleCBL-Teaching/Growth-Remodeling}
    \item \texttt{github.com/StanfordCBCL/svGrowth}, \texttt{svFSGe}
    \item \texttt{github.com/yale-humphrey-lab/FEMBE\_Plugin}, \texttt{FEFSG\_Plugin}
  \end{itemize}
  \vspace{0.6em}
  \textbf{Figure sources.} Result figures reproduced with the \texttt{gr} package;
  literature figures cited per-slide.
\end{frame}

